\section{Marco teórico}

\subsection{Principio de Arquímedes}

El principio de Arquímedes establece que un cuerpo sumergido total o parcialmente en un fluido experimenta una fuerza de empuje vertical hacia arriba, cuya magnitud es igual al peso del fluido desplazado \cite{serway2018,young2018}:

\begin{equation}
    E = \rho_f V_f g,
    \label{eq:empuje}
\end{equation}

donde $E$ es el empuje, $\rho_f$ es la densidad del fluido, $V_f$ es el volumen de fluido desplazado y $g$ es la aceleración de la gravedad.

En el experimento, una varilla se sumergió progresivamente en un beaker con agua colocado sobre una balanza. El aumento de la lectura de la balanza se atribuyó a la fuerza de reacción, igual y opuesta al empuje ejercido sobre la varilla, de acuerdo con la tercera ley de Newton.

Si la varilla posee un área transversal constante $A$, el volumen sumergido es $V_f=AL$. Por consiguiente, la diferencia de masa registrada por la balanza puede expresarse como

\begin{equation}
    \Delta m = (\rho_f A)L,
    \label{eq:deltam}
\end{equation}

donde $L$ es la longitud sumergida. La relación entre $\Delta m$ y $L$ es lineal, por lo que la pendiente $k$ del ajuste satisface $k=\rho_f A$. La densidad experimental del fluido se obtiene entonces mediante

\begin{equation}
    \rho_f = \frac{k}{A}.
    \label{eq:densidad}
\end{equation}

Para evaluar la consistencia del modelo, el intercepto del ajuste debe ser cercano a cero, pues un valor apreciable podría indicar la presencia de efectos sistemáticos o una lectura de referencia inadecuada.

\subsection{Ecuación de continuidad y teorema de Torricelli}

El vaciado de un recipiente con agua se describe mediante la ecuación de continuidad y el teorema de Torricelli. Para un fluido incompresible, la conservación del caudal entre el recipiente y el orificio de salida se expresa como

\begin{equation}
    A_1v_1=A_2v_2,
    \label{eq:continuidad}
\end{equation}

donde $A_1$ y $v_1$ son el área y la velocidad en el orificio, mientras que $A_2$ y $v_2$ corresponden al área y la velocidad de descenso del nivel en el recipiente. Esta relación representa la conservación de la masa del fluido \cite{young2018}.

El teorema de Torricelli, obtenido a partir de la ecuación de Bernoulli para un recipiente abierto a la atmósfera, permite calcular la velocidad de salida del agua por el orificio:

\begin{equation}
    v_1=\sqrt{2gh},
    \label{eq:torricelli}
\end{equation}

donde $h$ es la altura del nivel del agua medida desde el centro del orificio. Al combinar las ecuaciones de continuidad y de Torricelli, el tiempo teórico de vaciado puede escribirse en términos de áreas o diámetros como

\begin{equation}
    t_v=\left(\frac{A_2}{A_1}\right)\sqrt{\frac{2h}{g}}
    =\left(\frac{d_2}{d_1}\right)^2\sqrt{\frac{2h}{g}},
    \label{eq:tiempo_vaciado}
\end{equation}

donde $d_1$ es el diámetro del orificio y $d_2$ es el diámetro del recipiente. Esta expresión permite comparar el tiempo teórico de vaciado con el tiempo medido experimentalmente.

\subsection{Justificación del modelo experimental}

La determinación del empuje mediante la varilla y la balanza permite evaluar la relación entre el volumen de fluido desplazado y la lectura de la balanza. Debido a que el área transversal de la varilla se mantiene constante, se espera que la diferencia de masa aumente proporcionalmente con la longitud sumergida.

El experimento de vaciado permite estudiar la dependencia del tiempo necesario para vaciar el recipiente con respecto a la altura inicial del agua. Las mediciones de $h$ y $t$, junto con los diámetros $d_1$ y $d_2$, permiten comparar el comportamiento observado con el tiempo teórico obtenido mediante las ecuaciones de continuidad y de Torricelli.
